% ===== Block B06: 原文 p.47–54 | §3.1 共形块、§3.2 形式邻域 | 译手 tr-ch3a =====
\subsection{共形块}
% [原文 p.47]
本节中，对带形式邻域的稳定 $N$-点稳定曲线
\[
\mathfrak{X}=(C;Q_1,Q_2,\dots,Q_N;\eta_1,\eta_2,\dots,\eta_N)
\]
我们配备共形块或真空空间 $\cV^{\dagger}_{\vec{\lambda}}(\mathfrak{X})$ 及其对偶空间 $\cV_{\vec{\lambda}}(\mathfrak{X})$，其中 $\vec{\lambda}=(\lambda_1,\lambda_2,\dots,\lambda_N)$，$\lambda_j\in P_\ell$。为此我们首先定义广义仿射李代数 $\hat{\mathfrak{g}}_N$。

% [Tutor Note: 原文 "a stable $N$-pointed stable curve" 中 stable 重复出现，疑为笔误（应为 a stable $N$-pointed curve）。此处照原文译出，已在报告中登记 ERRATA。
%   核对：$N$-点稳定曲线的定义见 §1.2，与 "带形式邻域" 的附加数据 $\eta_j$ 无关。]

\begin{definition}[3.1]
设 $\mathfrak{g}$ 是复数域上的单李代数。李代数 $\hat{\mathfrak{g}}_N$ 定义为
\[
\hat{\mathfrak{g}}_N=\bigoplus_{j=1}^{N}\mathfrak{g}\otimes_{\C}\C((\xi_j))\oplus\C c
\]
并具有如下换位关系：
\begin{equation}
[(X_j\otimes f_j),(Y_j\otimes g_j)]=([X_j,Y_j]\otimes f_j g_j)+c\sum_{j=1}^{N}(X_j,Y_j)\operatorname{Res}_{\xi_j=0}(g_j\,\dd f_j) \tag{3.1}
\end{equation}
其中 $(a_j)$ 表示 $(a_1,a_2,\dots,a_N)$，而 $c$ 属于 $\hat{\mathfrak{g}}_N$ 的中心。
\end{definition}

% [Tutor Note: (3.1) 左端是"逐分量"的换位子：$X_j\otimes f_j$ 是第 $j$ 个分量的元素，当 $i\ne j$ 时 $[(X_i\otimes f_i),(Y_j\otimes g_j)]=0$；故右端只出现同一个 $j$ 的 $(X_j,Y_j)$。
%   核对归一化：$(X_j,Y_j)$ 是 $\mathfrak{g}$ 上的不变双线性型（通常取 Killing 型再除以 $2h^\vee$ 使长根平方长为 2）；中心项系数与这个归一化绑定，换归一化会改变 $c$ 的尺度。原文未说明是哪一种归一化。]

% [Tutor Note: (3.1) 的中心扩张项写成 $\operatorname{Res}(g_j\,\dd f_j)$ 而不是 $\operatorname{Res}(f_j\,\dd g_j)$，两者相差一个符号（$\dd(f_jg_j)$ 的留数为零）。抄写时勿把 $f,g$ 换位；下游的 KZ 方程系数与这个符号约定有关。]

% [Tutor Note: $\xi_j$ 是 $Q_j$ 处的局部坐标，$\C((\xi_j))$ 是形式 Laurent 级数域。"形式邻域"$\eta_j$ 的精确含义在 §3.2 才给出；此处的 $\mathfrak{X}$ 只是把曲线与局部坐标捆在一起的形式记号。]

我们还置
\begin{equation}
\hat{\mathfrak{g}}(\mathfrak{X})=\mathfrak{g}\otimes_{\C}H^0(C,\cO_C(*\sum_{j=1}^{N}Q_j)). \tag{3.2}
\end{equation}


由引理 1.21，我们有自然嵌入
\[
\iota:H^0(C,\cO_C(*\sum_{j=1}^{N}Q_j))\hookrightarrow\bigoplus_{j=1}^{N}\C((\xi_j)).
\]


以下我们常把 $H^0(C,\cO_C(*\sum_{j=1}^{N}Q_j))$ 看作 $\bigoplus_{j=1}^{N}\C((\xi_j))$ 的子空间。

由于闭黎曼曲面或结点曲线上亚纯一形式的留数之和为零，我们有如下引理。

% [Tutor Note: 结点曲线（nodal curve）即允许普通二重点的稳定曲线；留数定理在结点曲线上仍成立：在正规化后的每个不可约分支上分别用留数定理，结点处两个分支的留数相消。这是引理 3.2 能用 (3.1) 验证的关键。]

\begin{lemma}[3.2]
$\hat{\mathfrak{g}}(\mathfrak{X})$ 是 $\hat{\mathfrak{g}}_N$ 的李子代数。
\end{lemma}

% [Tutor Note: 引理 3.2 的验证只需看 (3.1) 的中心项：取 $f,g\in H^0(C,\cO_C(*\sum_jQ_j))$，则 $g\,\dd f$ 是亚纯一形式，其在 $Q_j$ 处的留数之和为零，故中心项消失；而 $fg$ 仍属于 $H^0$，故对换位子封闭。
%   注意这里"留数之和为零"是对所有 $j$ 求和，恰与 (3.1) 中对 $j$ 的求和匹配——这是定义中把中心项写成求和形式的动机。]

固定一个非负整数 $\ell$。对每个 $\vec{\lambda}=(\lambda_1,\dots,\lambda_N)\in(P_\ell)^N$，左 $\hat{\mathfrak{g}}_N$-模 $\cH_{\vec{\lambda}}$ 与右 $\hat{\mathfrak{g}}_N$-模 $\cH^{\dagger}_{\vec{\lambda}}$ 定义为
\[
\begin{aligned}
\cH_{\vec{\lambda}}&=\cH_{\lambda_1}\otimes_{\C}\cdots\otimes_{\C}\cH_{\lambda_N}\\
\cH^{\dagger}_{\vec{\lambda}}&=\cH^{\dagger}_{\lambda_1}\widehat{\otimes}_{\C}\cdots\widehat{\otimes}_{\C}\cH^{\dagger}_{\lambda_N}.
\end{aligned}
\]

% [Tutor Note: $P_\ell$ 是第 2 章引入的固定水平 $\ell$ 的整支配权（dominant integral weights）集合；$\lambda_j\in P_\ell$ 保证 $\cH_{\lambda_j}$ 是可积最高权模。$\ell$ 一旦固定，$\vec\lambda$ 的每个分量都取同一水平 $\ell$，这是共形块理论中的"水平守恒"前提。]

% [Tutor Note: 左模与右模的区别：$\cH^\dagger_{\vec\lambda}$ 是 $\cH_{\vec\lambda}$ 的对偶（带反号的作用），故 $g\in\hat{\mathfrak{g}}_N$ 从右边作用。第二行用"完备张量积"$\widehat{\otimes}_{\C}$（原文 \widehat{\otimes}）而非普通 $\otimes_{\C}$：对偶空间是拓扑对偶，需要完备化；第一行是普通张量积。抄写时勿把 $\otimes$ 与 $\widehat{\otimes}$ 混同。]

% [原文 p.48]
对每个元素 $X_j\in\mathfrak{g}$、$f(\xi_j)\in\C((\xi_j))$，$X_j[f_j]$ 在 $\cH_{\vec{\lambda}}$ 上的作用 $\rho_j$ 由
\begin{equation}
\rho_j(X_j[f_j])|v_1\otimes\cdots\otimes v_N\rangle
=|v_1\otimes\cdots\otimes v_{j-1}\otimes(X_j[f_j])v_j\otimes v_{j+1}\otimes\cdots\otimes v_N\rangle \tag{3.3}
\end{equation}
给出，其中 $|v_1\otimes\cdots\otimes v_N\rangle$ 表示 $|v_1\rangle\otimes\cdots\otimes|v_N\rangle$，$|v_j\rangle\in\cH_{\lambda_j}$。左 $\hat{\mathfrak{g}}_N$-作用由
\begin{equation}
(X_1\otimes f_1,\dots,X_N\otimes f_N)|v_1\otimes\cdots\otimes v_N\rangle
=\sum_{j=1}^{N}\rho_j(X_j[f_j])|v_1\otimes\cdots\otimes v_N\rangle \tag{3.4}
\end{equation}
给出。

% [Tutor Note: (3.3) 中 $X_j[f_j]$ 是 $X_j\otimes f_j$ 的简写，$\rho_j$ 是仿射李代数 $\hat{\mathfrak{g}}$（单个变量 $\xi_j$ 的那个）在 $\cH_{\lambda_j}$ 上的表示，不是 $\hat{\mathfrak{g}}_N$ 的表示；展开 (3.4) 时第 $j$ 项只作用在第 $j$ 个张量因子上，其余因子原样保留，各分量相加而不是组合成外积。
%   右作用 (3.5) 则是转置作用：$\langle u|(X\otimes f)$ 由 $\langle u|(X\otimes f)|v\rangle=\langle u|(X\otimes f)v\rangle$ 定义，这正是 (3.8) 的 $\hat{\mathfrak{g}}_N$-不变性 $\langle\Psi g|\Phi\rangle=\langle\Psi|g\Phi\rangle$ 的来源；若右作用另取符号，不变性会差一个负号。]

类似地，右 $\hat{\mathfrak{g}}_N$-作用在 $\cH^{\dagger}_{\vec{\lambda}}$ 上由
\begin{equation}
\langle u_1\otimes\cdots\otimes u_N|(X_1\otimes f_1,\dots,X_N\otimes f_N)
=\sum_{j=1}^{N}\langle u_1\otimes\cdots\otimes u_N|\rho_j(X_j[f_j]) \tag{3.5}
\end{equation}
定义。


作为李子代数，$\hat{\mathfrak{g}}(\mathfrak{X})$ 在 $\cH_{\vec{\lambda}}$ 与 $\cH^{\dagger}_{\vec{\lambda}}$ 上的作用为
\begin{align}
(X\otimes f)|v_1\otimes\cdots\otimes v_N\rangle
&=\sum_{j=1}^{N}\rho_j(X\otimes t_j(f))|v_1\otimes\cdots\otimes v_N\rangle \tag{3.6}\\
\langle u_1\otimes\cdots\otimes u_N|(X\otimes f)
&=\sum_{j=1}^{N}\langle u_1\otimes\cdots\otimes u_N|\rho_j(X\otimes t_j(f)). \tag{3.7}
\end{align}

% [Tutor Note: $t_j(f)$ 是 $f$ 在 $Q_j$ 处按局部参数 $\xi_j$ 的 Laurent 展开，即引理 1.21 的嵌入 $t$ 的第 $j$ 个分量。关键区别：在 (3.4) 中 $f_1,\dots,f_N$ 是各自独立的 Laurent 级数；而在 (3.6) 中 $t_1(f),\dots,t_N(f)$ 来自同一个整体亚纯函数 $f$，所以 (3.6) 只是 (3.4) 在子代数 $\hat{\mathfrak{g}}(\mathfrak{X})=\mathfrak{g}\otimes H^0(C,\cO_C(*\sum Q_j))$ 上的限制。
%   引理 1.21 给出 $t$ 是单射（需条件 (Q)），故 $\hat{\mathfrak{g}}(\mathfrak{X})$ 可无歧义地看作 $\hat{\mathfrak{g}}_N$ 的子代数。]

(2.24) 中的配对 $\langle\ |\ \rangle$ 诱导出一个完美双线性配对
\begin{align}
\langle\ |\ \rangle:\cH^{\dagger}_{\vec{\lambda}}\times\cH_{\vec{\lambda}}&\to\C \tag{3.8}\\
\big(\big(\langle u_1\otimes\cdots\otimes u_N|,|v_1\otimes\cdots\otimes v_N\rangle\big)\big)
&\mapsto\langle u_1|v_1\rangle\langle u_2|v_2\rangle\cdots\langle u_N|v_N\rangle, \nonumber
\end{align}
它是 $\hat{\mathfrak{g}}_N$-不变的：
\[
\langle\Psi(X_j\otimes f_j)|\Phi\rangle=\langle\Psi|(X_j\otimes f_j)\Phi\rangle.
\]



现在我们可以定义附着于 $\mathfrak{X}$ 的真空空间了。

\begin{definition}[3.3]
假设 $\mathfrak{X}$ 满足 §1.2 中的条件 (Q)。置
\begin{equation}
\cV_{\vec{\lambda}}(\mathfrak{X})=\cH_{\vec{\lambda}}/\hat{\mathfrak{g}}(\mathfrak{X})\cH_{\vec{\lambda}}. \tag{3.9}
\end{equation}
向量空间 $\cV_{\vec{\lambda}}(\mathfrak{X})$ 称为附着于 $\mathfrak{X}$ 的\emph{共真空空间}（space of covacua）。附着于 $\mathfrak{X}$ 的\emph{共形块}或\emph{真空空间}定义为
\begin{equation}
\cV^{\dagger}_{\vec{\lambda}}(\mathfrak{X})=\Hom_{\C}(\cV_{\vec{\lambda}}(\mathfrak{X}),\C). \tag{3.10}
\end{equation}
\end{definition}

% [Tutor Note: 条件 (Q) 是 §1.2 中的条件："曲线 $C$ 的每个不可约分支至少包含一个 $Q_j$"（见 p.18）。它的作用是保证引理 1.21 中 Laurent 展开映射 $t$ 为单射；否则整体亚纯函数可能在一个分支上恒为零而在别处非零，$t$ 有核，$\hat{\mathfrak{g}}(\mathfrak{X})$ 与 $\mathfrak{g}\otimes t(H^0)$ 的对应会出问题。
%   这也解释了下一段"添加标号为 0 的新点"：若 (Q) 不成立，就在缺点的分支上补一个标记点（权 0，即真空表示），使 (Q) 成立后再定义。]

当 $\mathfrak{X}$ 不满足性质 (Q) 时，我们用下面的定理 3.15 通过添加标号为 0 的新点来定义真空空间。

由定义 (3.10) 容易得到如下引理。

\begin{lemma}[3.4]
\begin{equation}
\cV^{\dagger}_{\vec{\lambda}}(\mathfrak{X})
=\{\,\langle\Psi|\in\cH^{\dagger}_{\vec{\lambda}}\mid\langle\Psi|\hat{\mathfrak{g}}(\mathfrak{X})=0\,\}. \tag{3.11}
\end{equation}
\end{lemma}

此外，配对 (3.8) 诱导出一个完美配对
\begin{equation}
\langle\ |\ \rangle:\cV^{\dagger}_{\vec{\lambda}}(\mathfrak{X})\times\cV_{\vec{\lambda}}(\mathfrak{X})\to\C. \tag{3.12}
\end{equation}

% [Tutor Note: 核对 (3.11)：$\cV_{\vec\lambda}=\cH_{\vec\lambda}/\hat{\mathfrak{g}}\cH_{\vec\lambda}$ 的对偶是 $\hat{\mathfrak{g}}\cH_{\vec\lambda}$ 的零化子，即所有满足 $\langle\Psi|gv\rangle=0$（对一切 $g\in\hat{\mathfrak{g}}(\mathfrak{X})$、$v\in\cH_{\vec\lambda}$）的 $\langle\Psi|$；用 (3.8) 的不变性可把它写成 $\langle\Psi|g=0$。
%   (3.12) 的完美性是把 (3.8) 限制到零化子与商空间上得到的，这是标准线性代数，但前提是 $\cV_{\vec\lambda}$ 的维数有限——这正是定理 3.5 要证明的。
%   定理 3.5 是本章技术核心之一：$\cV^\dagger_{\vec\lambda}$ 定义为 $\Hom_{\C}(\cV_{\vec\lambda},\C)$，在不知道 $\cV_{\vec\lambda}$ 有限维时这只是代数对偶；定理 3.5 同时给出两个空间有限维，才使 (3.12) 的"完美配对"名副其实。证明在 p.49–52（命题 3.6、3.7 与引理 3.8–3.10）：用 $\hat{\mathfrak{g}}_N$ 的滤过与 $\cH_{\vec\lambda}$ 的可积性把商空间压到有限维。]

按我们的定义，$\cV^{\dagger}_{\vec{\lambda}}(\mathfrak{X})$ 是否为有限维向量空间并不清楚。现在我们证明如下重要事实。

\begin{theorem}[3.5]
$\cV_{\vec{\lambda}}(\mathfrak{X})$ 与 $\cV^{\dagger}_{\vec{\lambda}}(\mathfrak{X})$ 都是有限维向量空间。
\end{theorem}

% [原文 p.49]
在第 IV 章我们将证明，$\dim_{\C}\cV^{\dagger}_{\vec{\lambda}}(\mathfrak{X})$ 只依赖于曲线 $C$ 的亏格 $g(C)$ 与 $\vec{\lambda}$。

在 [TUY] 中，定理 3.5 是用 Gabber 定理证明的。这里我们给出 Suzuki [Su] 的一个初等证明。

首先我们引入 $\hat{\mathfrak{g}}_N$ 的李子代数 $\mathfrak{g}[1]$ 与 $\hat{\mathfrak{g}}_N^{(m)}$：
\begin{align}
\mathfrak{g}[1]&=\{\,(X\otimes 1,\dots,X\otimes 1)\in\hat{\mathfrak{g}}\mid X\in\mathfrak{g}\,\} \tag{3.13}\\
\hat{\mathfrak{g}}_N^{(m)}&=\mathfrak{g}\otimes_{\C}\bigoplus_{j=1}^{N}\C[\xi_j^{-1}]\xi_j^{-m},\qquad m=0,1,\dots. \tag{3.14}
\end{align}

% [Tutor Note: $\mathfrak{g}[1]$ 是"常值环路"构成的子代数（每个分量取同一个 $X\otimes1$）；$\hat{\mathfrak{g}}_N^{(m)}$ 是 $\mathfrak{g}\otimes\C[\xi_j^{-1}]\xi_j^{-m}$，即 $\xi_j^{-m},\xi_j^{-m-1},\dots$ 张成的"深负模"部分。
%   核对它们构成李子代数：$[X\otimes1,Y\otimes1]=[X,Y]\otimes1$ 仍在 $\mathfrak{g}[1]$ 中；$[X\otimes1,Y\otimes\xi^{-k}]=[X,Y]\otimes\xi^{-k}$ 且中心项因 $\dd1=0$ 为零，故仍在 $\hat{\mathfrak{g}}_N^{(m)}$ 中。]
%   注意 (3.13) 原文把 $\hat{\mathfrak{g}}$ 写成不带下标 $N$ 的形式，但集合里的元素是长度 $N$ 的多元组，按上下文应为 $\hat{\mathfrak{g}}_N$；已照原文排版并在报告 ERRATA 中说明。]

注意对任意正整数 $m$，$\mathfrak{g}[1]\oplus\hat{\mathfrak{g}}_N^{(m)}$ 也是 $\hat{\mathfrak{g}}_N$ 的李子代数。下面我们证明如下命题。

\begin{proposition}[3.6]
若向量空间
\begin{equation}
W_m=\cH_{\vec{\lambda}}/(\mathfrak{g}[1]\oplus\hat{\mathfrak{g}}_N^{(m)})\cH_{\vec{\lambda}} \tag{3.15}
\end{equation}
对某个正整数 $m$ 是有限维的，则定理 3.5 成立。
\end{proposition}

\noindent\textit{证明.}\quad 我们引入 $\mathfrak{g}$ 的一个滤过 $\{F_\bullet\}$：
\begin{equation}
F_p\mathfrak{g}=
\begin{cases}
\mathfrak{g}\otimes_{\C}\bigoplus_{j=1}^{N}\C[\xi_j^{-1}]\xi_j^{-p}\oplus\C\cdot c, & p\ge0,\\[2pt]
\mathfrak{g}\otimes_{\C}\bigoplus_{j=1}^{N}\C[\xi_j^{-1}]\xi_j^{-p}, & p<0.
\end{cases} \tag{3.16}
\end{equation}
于是 $\cH_{\vec{\lambda}}$ 有一个由 $\hat{\mathfrak{g}}_N$ 的滤过诱导的自然滤过，即
\begin{equation}
F_p\cH_{\vec{\lambda}}=F_pU(\hat{\mathfrak{g}}_N)\cdot V_{\vec{\lambda}}, \tag{3.17}
\end{equation}
其中 $U(\hat{\mathfrak{g}}_N)$ 表示李代数 $\hat{\mathfrak{g}}_N$ 的包络代数，而滤过由
\[
F_pU(\hat{\mathfrak{g}}_N)=\sum_{p_1+\cdots+p_N\le p}F_{p_1}\hat{\mathfrak{g}}_N\otimes\cdots\otimes F_{p_N}\hat{\mathfrak{g}}_N
\]
引入；并且
\[
V_{\vec{\lambda}}=V_{\lambda_1}\otimes\cdots\otimes V_{\lambda_N}.
\]

% [Tutor Note: (3.16) 的关键在 $p\ge0$ 时多出的中心项 $\C\cdot c$：中心元 $c$ 被放在滤过次数 $0$，因此 $F_0\mathfrak{g}$ 已经包含 $c$。这保证了 $F_p\mathfrak{g}$ 对换位子封闭（(3.1) 的中心项是 $c$ 的倍数）。
%   核对：$[F_p\mathfrak{g},F_q\mathfrak{g}]\subset F_{p+q}\mathfrak{g}$ 对 $p,q$ 的符号都成立，因为 $\xi^{-p}\xi^{-q}=\xi^{-(p+q)}$ 且 $\operatorname{Res}(g\,\dd f)$ 只涉及 $\xi^{-1}$ 系数。]

由定理 2.4，$\cH_{\vec{\lambda}}$ 作为 $U(\mathfrak{g}\otimes_{\C}\bigoplus_{j=1}^{N}\C[\xi_j^{-1}])$-模由 $V_{\vec{\lambda}}$ 生成，其中 $U(\mathfrak{q})$ 表示李代数 $\mathfrak{q}$ 的包络代数。于是我们有
\[
F_pU(\hat{\mathfrak{g}}_N)\cdot F_q\cH_{\vec{\lambda}}\subset F_{p+q}\cH_{\vec{\lambda}}.
\]
注意我们还有
\[
F_p\cH_{\vec{\lambda}}=\{0\}\qquad(p<0).
\]

让我们考虑 $\hat{\mathfrak{g}}_N$ 的子代数上的诱导滤过以及 $\cV_{\vec{\lambda}}(\mathfrak{X})$ 与 $W_m$ 上的商滤过，并用 $\operatorname{gr}^F_{\bullet}(\ )$ 表示相应的分次对象。于是我们有如下分次 $\C$-向量空间的短正合列：
\[
\begin{aligned}
0&\longrightarrow\operatorname{gr}^F_{\bullet}(\hat{\mathfrak{g}}(\mathfrak{X}))\longrightarrow\operatorname{gr}^F_{\bullet}\cH_{\vec{\lambda}}\longrightarrow\operatorname{gr}^F_{\bullet}\cV_{\vec{\lambda}}(\mathfrak{X})\longrightarrow0\\
0&\longrightarrow\operatorname{gr}^F_{\bullet}((\hat{\mathfrak{g}}_N^{(m)}\oplus\mathfrak{g}[1])\cH_{\vec{\lambda}})\longrightarrow\operatorname{gr}^F_{\bullet}\cH_{\vec{\lambda}}\longrightarrow\operatorname{gr}^F_{\bullet}W_m\longrightarrow0
\end{aligned}
\]
由 Riemann-Roch 定理，对任意整数 $m\ge 2g-1$ 我们有
\begin{equation}
\operatorname{gr}^F_{\bullet}((\hat{\mathfrak{g}}_N^{(m)}\oplus\mathfrak{g}[1])\cH_{\vec{\lambda}})
\subset\operatorname{gr}^F_{\bullet}(\hat{\mathfrak{g}}(\mathfrak{X})\cH_{\vec{\lambda}}) \tag{3.18}
\end{equation}

% [Tutor Note: 证明策略：由 (3.18)，$(\hat{\mathfrak{g}}_N^{(m)}\oplus\mathfrak{g}[1])\cH_{\vec\lambda}\subset\hat{\mathfrak{g}}(\mathfrak{X})\cH_{\vec\lambda}$（在分次层面包含）；于是有满射 $W_m\twoheadrightarrow\cV_{\vec\lambda}(\mathfrak{X})$，从而 $\dim\cV_{\vec\lambda}\le\dim W_m<\infty$，再取对偶得定理 3.5。
%   这是全证明的技术核心。Riemann-Roch 在这里用来保证 $m\ge2g-1$ 时，截断子代数 $\hat{\mathfrak{g}}_N^{(m)}$ 的深负模部分可以在分次层面被 $\hat{\mathfrak{g}}(\mathfrak{X})$ 中整体亚纯函数的 Laurent 展开"模拟"；原文没有写出 $h^0(C,\cO_C(m\sum Q_j))$ 的具体计算（$m\ge2g-1$ 时 $h^1=0$），这是要补的一步，★★★。]

% [原文 p.50]
这两者作为 $\operatorname{gr}^F_{\bullet}\cH_{\vec{\lambda}}$ 的子空间。取定这样的 $m$。于是由 (3.18) 及上述正合列，我们得到满同态
\[
\operatorname{gr}^F_{\bullet}W_m\longrightarrow\operatorname{gr}^F_{\bullet}\cV_{\vec{\lambda}}(\mathfrak{X})\longrightarrow0.
\]
因此 $\operatorname{gr}^F_{\bullet}\cV_{\vec{\lambda}}(\mathfrak{X})$ 是有限维的。这就推出 $\cV_{\vec{\lambda}}(\mathfrak{X})$ 的有限维性。\hfill$\square$

% [Tutor Note: 原文 p.49 以 (3.18) 结束，p.50 直接以 "as subspaces of gr^F H" 开头，句子前半（主语与谓语）在换页处缺失，疑为原书排版脱漏。按上下文，主语应为上一段所说的两个分次空间 $\operatorname{gr}^F_\bullet\cV_{\vec\lambda}(\mathfrak{X})$ 与 $\operatorname{gr}^F_\bullet W_m$；本译以「这两者」承接，并在报告中登记 ERRATA。]

\begin{proposition}[3.7]
对任意正整数 $m\ge 2g-1$，由 (3.15) 定义的向量空间
\[
W_m=\cH_{\vec{\lambda}}/(\mathfrak{g}[1]\oplus\hat{\mathfrak{g}}_N^{(m)})\cH_{\vec{\lambda}}
\]
是有限维的。
\end{proposition}

\noindent\textit{证明.}\quad 令 $W=W_m$，
\[
\bar{\mathfrak{g}}=\bigoplus_{j=1}^{N}\Bigl(\mathfrak{g}\otimes_{\C}\Bigl(\bigoplus_{n=0}^{m-1}\C\xi_j^{-n}\Bigr)\Bigr).
\]
于是我们有同构
\begin{equation}
\bar{\mathfrak{g}}\cong\hat{\mathfrak{g}}_N^{(0)}/\hat{\mathfrak{g}}_N^{(m)}. \tag{3.19}
\end{equation}
由于 $\hat{\mathfrak{g}}_N^{(0)}$ 是 $\hat{\mathfrak{g}}_N$ 的李子代数，而 $\hat{\mathfrak{g}}_N^{(m)}$ 是 $\hat{\mathfrak{g}}_N^{(0)}$ 的理想，由上述同构我们可以把 $\bar{\mathfrak{g}}$ 看作一个有限维李代数。于是向量空间
\[
U=\cH_{\vec{\lambda}}/\hat{\mathfrak{g}}_N^{(m)}\cH_{\vec{\lambda}}
\]
可以看作 $U(\bar{\mathfrak{g}})$-模。设 $\bar V$ 表示 $V_{\vec{\lambda}}$ 在 $U$ 中的像。则 $\bar V$ 是有限维向量空间，并且有
\begin{equation}
U=U(\bar{\mathfrak{g}})\cdot\bar V. \tag{3.20}
\end{equation}

我们在 $U(\hat{\mathfrak{g}}_N^{(0)})$ 与 $\cH_{\vec{\lambda}}$ 上如下引入滤过 $G_{\bullet}$：
\[
\begin{aligned}
G_kU(\hat{\mathfrak{g}}_N^{(0)})&=
\begin{cases}
\{0\}, & k<0,\\
\C\cdot1, & k=0,\\
G_{k-1}U(\hat{\mathfrak{g}}_N^{(0)})+\hat{\mathfrak{g}}_N^{(0)}G_{k-1}U(\hat{\mathfrak{g}}_N^{(0)}), & k>0,
\end{cases}\\
G_k\cH_{\vec{\lambda}}&=G_kU(\hat{\mathfrak{g}}_N^{(0)})\cdot V_{\vec{\lambda}}.
\end{aligned}
\]
于是容易证明
\[
\begin{aligned}
[G_kU(\hat{\mathfrak{g}}_N^{(0)}),G_lU(\hat{\mathfrak{g}}_N^{(0)})]&\subset G_{k+l-1}U(\hat{\mathfrak{g}}_N^{(0)}),\\
G_kU(\hat{\mathfrak{g}}_N^{(0)})G_k\cH_{\vec{\lambda}}&\subset G_{k+l}\cH_{\vec{\lambda}}.
\end{aligned}
\]

% [Tutor Note: 上面第二式原文左右两侧的 $U$ 与 $\cH$ 前都印作 $G_k$，与右端 $G_{k+l}$ 不匹配；按第一式的模式（以及左边应为 $G_k$、$G_l$ 两个因子），第二个因子应为 $G_l$，即 $G_kU\cdot G_l\cH\subset G_{k+l}\cH$。已照原文排版，并在报告 ERRATA 中登记。]

在 $U(\bar{\mathfrak{g}})$ 上，我们利用同构 (3.19) 引入商滤过。类似地，由 $\cH_{\vec{\lambda}}$ 的滤过在 $W$ 与 $U$ 上引入商滤过。于是我们有分次代数的同构
\[
\operatorname{gr}^G_{\bullet}U(\bar{\mathfrak{g}})\cong S(\bar{\mathfrak{g}}),
\]
其中 $S(\bar{\mathfrak{g}})$ 是向量空间 $\bar{\mathfrak{g}}$ 的对称代数。由 (3.20)，我们有分次向量空间的同构
\[
\operatorname{gr}^G_{\bullet}U\cong S(\bar{\mathfrak{g}})\cdot\bar V.
\]

我们用 $S_k(\bar{\mathfrak{g}})$ 记 $S(\bar{\mathfrak{g}})$ 的 $k$ 次部分。现在为证明定理 3.5，首先注意容易证明下述引理。

% [原文 p.51]
\begin{lemma}[3.8]
对每个整数 $k$，下述图表交换，且所有态射都是满射。
\[
\begin{array}{ccc}
G_kU & \xrightarrow{\ r_k\ } & S_k(\bar{\mathfrak{g}})\cdot\bar V\\[6pt]
\big\downarrow & & \big\downarrow{\scriptstyle\varphi_k}\\[6pt]
G_kW & \longrightarrow & \operatorname{gr}^G_k W.
\end{array}
\]
\end{lemma}

由于 $\bar{\mathfrak{g}}$ 与 $V$ 都是有限维的，上述命题 3.7 推出 $\operatorname{gr}^G_k W$ 是有限维的。因此，为证明定理 3.5，只需证明下述引理。

\begin{lemma}[3.9]
存在整数 $K$，使得对任意 $k\ge K$ 都有
\begin{equation}
\operatorname{gr}^G_k W=0. \tag{3.21}
\end{equation}
\end{lemma}

注意由滤过 $G_{\bullet}$ 的定义，若 (3.21) 对某个整数 $k=K$ 成立，则 (3.21) 对任意 $k>K$ 也成立。因此我们只需对某个 $k$ 证明 (3.21)。现在取充分大的 $k$，并令
\[
\pi_k=\varphi_k\circ r_k:G_kU\longrightarrow\operatorname{gr}^G_k W.
\]

\begin{lemma}[3.10]
对任意元素 $X\in\mathfrak{g}$，存在正整数 $L$，使得对任意 $l\ge L$、$j=1,\dots,N$、$n=0,1,\dots,m-1$ 以及 $|u\rangle\in G_{k-l}U$，都有
\[
\pi_k\bigl((\rho_j(X\otimes\xi_j^{-n})^l)|u\rangle\bigr)=0.
\]
\end{lemma}

由于 $S(\bar{\mathfrak{g}})$ 的任意元素都是形如 $(X\otimes\xi_j^{-n})^l$ 的乘积的有限线性组合，引理 3.10 蕴涵引理 3.9。

\noindent\textit{引理 3.10 的证明.}\quad 固定 $j$、$n$ 与 $\alpha\in\Delta_+$，并令
\[
E=X_{\alpha},\qquad F=X_{-\alpha},\qquad H=[X_{\alpha},X_{-\alpha}].
\]
由于 $E\otimes\xi_j^{-n}$ 与 $F\otimes\xi_j^{-n}$ 在 $\cH_{\vec{\lambda}}$ 上幂零地作用，而 $S_{k-l}(\bar{\mathfrak{g}})\cdot\bar V$ 对 $l=0,1,\dots,k$ 是有限维的，我们可以找到一个正整数 $L$，使得
\begin{equation}
\pi_k\bigl((\rho_j(E\otimes\xi_j^{-n})^l)|u\rangle\bigr)=0,\qquad
\pi_k\bigl((\rho_j(F\otimes\xi_j^{-n})^l)|u\rangle\bigr)=0 \tag{3.22}
\end{equation}
对任意 $l\ge L$、$j=1,\dots,N$、$n=0,1,\dots,m-1$ 以及 $|u\rangle\in G_{k-l}U$ 成立。因此，在这些情形下引理成立。

现在对任意元素 $|u\rangle\in G_{k-s-t}U$，令
\[
|v\rangle=\rho_j\bigl((H\otimes\xi_j^{-n})^{s-1}(F\otimes\xi_j^{-n})^{t+1}\bigr)|u\rangle,
\]
其中我们假设 $s\ge1$ 且 $t\ge0$。再令
\[
E[1]=\bigoplus_{i=1}^{N}E\otimes\xi_i^{0}.
\]
于是我们有
\begin{equation}
\begin{aligned}
E[1]|v\rangle=\sum_{i=1}^{N}\rho_i(E\otimes1)|v\rangle
={}&\bigl[\rho_j(E\otimes1),\rho_j\bigl((H\otimes\xi_j^{-n})^{s-1}(F\otimes\xi_j^{-n})^{t+1}\bigr)\bigr]|u\rangle\\
&+\rho_j\bigl((H\otimes\xi_j^{-n})^{s-1}(F\otimes\xi_j^{-n})^{t+1}\bigr)E[1]|u\rangle.
\end{aligned} \tag{3.23}
\end{equation}
注意
\[
E[1]|u\rangle\in G_{k-s-t}U.
\]

% [Tutor Note: 这里 $E=X_\alpha$、$F=X_{-\alpha}$、$H=[X_\alpha,X_{-\alpha}]$ 构成 $\mathfrak{sl}_2$ 三元组；$E\otimes\xi_j^{-n}$ 与 $F\otimes\xi_j^{-n}$ 的幂零性来自定理 2.5（可积模的根向量局部幂零作用）。证明用 $H$ 把 $F$ 的作用「提升」到次数更高处，再由 (3.22) 的两种情形（$E$ 型与 $F$ 型）归纳，这是证明有限维性的标准手法。]

% [原文 p.52]
而由 $W$ 的定义 (3.15)，我们有
\[
\pi_k(E[1]|v\rangle)=0.
\]
于是由上述等式 (3.23) 推得
\[
\begin{aligned}
0=\pi_k(E[1]|v\rangle)={}&-2(s-1)\pi_k\bigl(\rho_j((H\otimes\xi_j^{-n})^{s-2}(F\otimes\xi_j^{-n})^{t+1}E\otimes\xi_j^{-n})|u\rangle\bigr)\\
&+(t+1)\pi_k\bigl(\rho_j((H\otimes\xi_j^{-n})^{s}(F\otimes\xi_j^{-n})^{t})|u\rangle\bigr)\\
&+\pi_k\bigl(\rho_j((H\otimes\xi_j^{-n})^{s-1}(F\otimes\xi_j^{-n})^{t+1})E[1]|u\rangle\bigr).
\end{aligned}
\]
由此得到
\[
\begin{aligned}
&\pi_k\bigl(\rho_j((H\otimes\xi_j^{-n})^{s}(F\otimes\xi_j^{-n})^{t})|u\rangle\bigr)\\
&\quad=\frac{1}{t+1}\Bigl\{2(s-1)\pi_k\bigl(\rho_j((H\otimes\xi_j^{-n})^{s-2}(F\otimes\xi_j^{-n})^{t+1}E\otimes\xi_j^{-n})|u\rangle\bigr)\\
&\qquad+\pi_k\bigl(\rho_j((H\otimes\xi_j^{-n})^{s-1}(F\otimes\xi_j^{-n})^{t+1})|u\rangle\bigr)\Bigr\}.
\end{aligned}
\]
这一等式表明，我们可以通过减少 $s$ 来增大 $t$，使得最终项 $\rho_j(H\otimes\xi_j^{-n})$ 消失，而右端只出现 $\rho_j(E\otimes\xi_j^{-n})$、$\rho_j(F\otimes\xi_j^{-n})$ 与 $E[1]$ 的项。因此由 (3.22)，当 $l\ge L$ 时我们也有
\[
\pi_k(\rho_j(H\otimes\xi_j^{-n})^{l}|u\rangle)=0.
\]
现在，由于李代数 $\mathfrak{g}$ 的任意元素都是 $X_{\alpha}$、$X_{-\alpha}$、$[X_{\alpha},X_{-\alpha}]$（$\alpha\in\Delta_+$）的线性组合，我们便得到所要的结论。\hfill$\square$

于是定理 3.5 得证。

% [Tutor Note: 上面「由此得到」的公式按原文照排：第二个花括号内的末项原文写作 $\pi_k(\rho_j(\cdots)|u\rangle)$，但上一式中对应的项是 $\pi_k(\rho_j(\cdots)E[1]|u\rangle)$，两处相差一个 $E[1]$。按推导，该处应保留 $E[1]$ 才与 (3.23) 一致；已照原文排版并在报告 ERRATA 中登记。]

\subsection{形式邻域}

我们令 $\cD$ 为形式 Laurent 级数域 $\C((\xi))$ 作为 $\C$-代数的自同构群 $\Aut\C((\xi))$。群 $\cD$ 也可以看作形式幂级数环 $\C[[\xi]]$ 的自同构群。存在自然同构
\[
\cD\cong\Bigl\{\sum_{n=0}^{\infty}a_n\xi^{n+1}\ \Bigm|\ a_0\ne0\Bigr\},\qquad h\mapsto h(\xi),
\]
其中 $h,g\in\cD$ 的复合 $h\circ g$ 对应于形式幂级数 $h(g(\xi))$。以下我们把 $\cD$ 与 $\{\sum_{n=0}^{\infty}a_n\xi^{n+1}\mid a_0\ne0\}$ 等同起来。

对正整数 $p$，令
\[
\cD^{p}=\{h\in\cD\mid h(\xi)=\xi+a_p\xi^{p+1}+\cdots\}.
\]
于是我们有滤过
\[
\cD=\cD^{0}\supset\cD^{1}\supset\cD^{2}\supset\cdots.
\]
再令
\[
\underline{\mathfrak{d}}=\C[[\xi]]\xi\frac{\dd}{\dd\xi},\qquad
\underline{\mathfrak{d}}^{p}=\C[[\xi]]\xi^{p+1}\frac{\dd}{\dd\xi},\quad p=0,1,2,\dots
\]
于是我们有滤过
\[
\underline{\mathfrak{d}}=\underline{\mathfrak{d}}^{0}\supset\underline{\mathfrak{d}}^{1}\supset\underline{\mathfrak{d}}^{2}\supset\cdots.
\]


% [原文 p.53]
对任意 $\underline l\in\underline{\mathfrak d}$ 与 $f(\xi)\in\C[[\xi]]$，定义 $\exp(\underline l)(f(\xi))$ 为
\[
\exp(\underline l)(f(\xi))=\sum_{k=0}^{\infty}\frac{1}{k!}(\underline l^k f(\xi)).
\]
再令
\[
\cD_+^{0}=\{h\in\cD\mid h(\xi)=a\xi+a_1\xi^{2}+\cdots,\quad a>0\},
\]
\[
\underline{\mathfrak d}_+^{0}=\Bigl\{l(\xi)\frac{\dd}{\dd\xi}\ \Bigm|\ l(\xi)=\alpha\xi+\alpha_1\xi^{2}+\cdots,\quad\alpha\in\R\Bigr\}.
\]
于是我们有下述结果。

\begin{lemma}[3.11]
指数映射
\[
\begin{array}{rcl}
\exp:\underline{\mathfrak d}&\longrightarrow&\cD\\
\underline l&\longmapsto&\exp(\underline l)
\end{array}
\]
是满射。此外，指数映射诱导同构
\[
\exp:\underline{\mathfrak d}_+^{0}\cong\cD_+^{0}.
\]
由于对任意整数 $n$ 都有
\[
\exp\Bigl(2\pi n\sqrt{-1}\xi\frac{\dd}{\dd\xi}\Bigr)=\id,
\]
故指数映射在 $\underline{\mathfrak d}$ 上不是单射。
\end{lemma}

对能动张量 $T(z)$ 与任意元素 $\underline l\in\cD_+^{0}$，我们定义 $\exp(T[\underline l])$ 为
\[
\exp(T[\underline l])=\sum_{k=0}^{\infty}\frac{1}{k!}T[\underline l]^{k}.
\]
于是 $\exp(T[\underline l])$ 从左作用于 $\cH_{\lambda}$，从右作用于 $\cH_{\lambda}^{\dagger}$。

由引理 3.11，对任意自同构 $h\in\cD_+^{0}$，存在唯一的 $\underline l\in\underline{\mathfrak d}_+^{0}$ 使得 $\exp(\underline l)=h$。现在对 $h\in\cD_+^{0}$ 定义算子 $G[h]$ 为
\[
G[h]=\exp(-T[\underline l]),
\]
其中 $\exp(\underline l)=h$。于是 $G[h]$ 从左作用于 $\cH_{\lambda}$，从右作用于 $\cH_{\lambda}^{\dagger}$。因此我们有下述重要结果。

\begin{theorem}[3.12]
对任意 $h\in\cD_+^{0}$、$f(\xi)\dd\xi\in\C((\xi))\dd\xi$、$g(\xi)\in\C((\xi))$ 以及 $\underline l=l(\xi)\frac{\dd}{\dd\xi}\in\C((\xi))\frac{\dd}{\dd\xi}$，作为 $\cH_{\lambda}$ 与 $\cH_{\lambda}^{\dagger}$ 上的算子，我们有下述等式：
\begin{enumerate}
\item[(1)] $G[h](X\otimes f(\xi))G[h]^{-1}=X\otimes f(h(\xi))$，\quad $f(\xi)\in\C((\xi))$。
\item[(2)] $G[h_1\circ h_2]=G[h_1]G[h_2]$，\quad 对 $h_1,h_2\in\cD^{p}$，$p\ge1$。
\item[(3)] $G[h]T[\underline l]G[h]^{-1}=T[\ad(h)(\underline l)]+\frac{c_v}{12}\operatorname{Res}_{\xi=0}\bigl(\{h(\xi);\xi\}l(\xi)\dd\xi\bigr)$。
\end{enumerate}
其中 $\underline l=l(\xi)\frac{\dd}{\dd\xi}$，而 $\{f(\xi);\xi\}$ 是 Schwarz 导数。
\end{theorem}

其证明可由直接计算容易地给出。由这个定理我们容易推出下述命题。

\begin{proposition}[3.13]
对任意 $h_j\in\cD_+^{0}$（$j=1,2,\dots,N$）以及带形式坐标的 $N$-点曲线
\[
\mathfrak{X}=(C;Q_1,Q_2,\dots,Q_N;\xi_1,\xi_2,\dots,\xi_N),
\]
令
\[
\mathfrak{X}_{(h)}=(C;Q_1,Q_2,\dots,Q_N;h_1(\xi_1),h_2(\xi_2),\dots,h_N(\xi_N)).
\]

% [原文 p.54]
则同构 $G[h_1]\widehat{\otimes}\cdots\widehat{\otimes}G[h_N]$
\[
\begin{aligned}
&\cH_{\vec{\lambda}}^{\dagger}=\cH_{\lambda_1}^{\dagger}\widehat{\otimes}\cdots\widehat{\otimes}\cH_{\lambda_N}^{\dagger}
\longrightarrow\cH_{\vec{\lambda}}^{\dagger}=\cH_{\lambda_1}^{\dagger}\widehat{\otimes}\cdots\widehat{\otimes}\cH_{\lambda_N}^{\dagger}\\
&\langle\phi_1\widehat{\otimes}\phi_2\widehat{\otimes}\cdots\widehat{\otimes}\phi_N|
\longmapsto\langle\phi_1G[h_1]\widehat{\otimes}\phi_2G[h_2]\widehat{\otimes}\cdots\widehat{\otimes}\phi_NG[h_N]|
\end{aligned}
\]
诱导出典范同构
\[
G[h_1]\widehat{\otimes}\cdots\widehat{\otimes}G[h_N]:\cV_{\vec{\lambda}}^{\dagger}(\mathfrak{X})\longrightarrow\cV_{\vec{\lambda}}^{\dagger}(\mathfrak{X}_{(h)}).
\]
\end{proposition}

\begin{corollary}[3.14]
若两条带形式坐标的稳定 $N$-点曲线
\[
\begin{aligned}
\mathfrak{X}&=(C;Q_1,Q_2,\dots,Q_N;\xi_1,\xi_2,\dots,\xi_N),\\
\mathfrak{X}'&=(C;Q_1,Q_2,\dots,Q_N;\eta_1,\eta_2,\dots,\eta_N)
\end{aligned}
\]
诱导出相同的带一阶邻域的稳定 $N$-点曲线，则存在典范同构
\[
\cV_{\vec{\lambda}}^{\dagger}(\mathfrak{X})\sim\cV_{\vec{\lambda}}^{\dagger}(\mathfrak{X}').
\]
\end{corollary}

\noindent\textit{证明.}\quad 由假设，存在 $h_j\in\cD^{1}$ 使得 $\eta_j=h_j(\xi_j)$，$j=1,\dots,N$。于是由上述定理即得所要结论。\hfill$\square$

% [Tutor Note: 命题 3.13 中的曲线记作 $\mathfrak{X}=(C;Q_1,\dots,Q_N;\xi_1,\dots,\xi_N)$（形式坐标），而 §3.1 开头把同一对象写成 $(C;Q_1,\dots,Q_N;\eta_1,\dots,\eta_N)$（形式邻域）；两条曲线的差别只在最后一项的记号。推论 3.14 里两条曲线分别用 $\xi$ 与 $\eta$ 表示坐标，正好把这两种写法并列，可据此确认二者是同一种数据的两种记法。
%   这正是「共形块对形式坐标的依赖是良定义的」的确切含义：改变坐标 $h_j$ 时，$G[h_1]\widehat\otimes\cdots\widehat\otimes G[h_N]$ 给出 $\cV^\dagger_{\vec\lambda}(\mathfrak{X})\to\cV^\dagger_{\vec\lambda}(\mathfrak{X}_{(h)})$ 的典范同构；而推论 3.14 说，一阶邻域相同（即 $\eta_j=h_j(\xi_j)$ 且 $h_j\in\cD^1$）的两组坐标给出同构的共形块。因此 $\cV^\dagger_{\vec\lambda}$ 只依赖「曲线 + 各点处形式坐标的一阶邻域」，不依赖坐标的高阶项。]

